%Dokumentenklasse
\documentclass[
	a4paper,										%Format
	12pt,											%Schriftgröße
	twoside,										%Zweiseitig
	headsepline,								%horizontale Linie oben
	footsepline,									%horizontale Linie unten
	captions=tableheading,					%Überschrift über Tabelle (nicht unter)
]{scrartcl}

\renewcommand{\thesubsection}{\thesection.\arabic{subsection}.}
\renewcommand{\thesubsubsection}{\thesection.\arabic{subsection}.\arabic{subsubsection}. \hspace{2pt}}
\renewcommand{\theparagraph}{\thesection.\arabic{subsection}.\arabic{subsubsection}.\arabic{paragraph}.}

%Pakete:
\usepackage[exponent-product = \cdot, output-decimal-marker = {,}]{siunitx}
\usepackage{textalpha}
\usepackage{floatrow}
\floatsetup[table]{style=plaintop}
\floatsetup[longtable]{LTcapwidth=table}
%\usepackage{overcite}
\usepackage{siunitx}
\usepackage[ngerman]{babel} 			%Neue deutsche Rechtschreibung
\usepackage{grffile}
\usepackage{graphicx} 					%Bilder einbinden
\usepackage{xcolor}						%Farben für Schrift und Plots
\usepackage{pgfplots}
\pgfplotsset{compat=1.18}					%Plots und Tabellen aus Datensätzen.
\usepackage{amsmath}						%Mathe
\usepackage{amssymb}					%Mathesymbole
\usepackage{longtable}					%Seitenübergreifende Tabellen zulassen
\usepackage{float}
\usepackage{tabularx}
\usepackage{geometry}					%Seitenränder
\usepackage{ae}								%skalierbare Schriften (im PDF-File)
\usepackage[activate]{microtype}		%optischer Randausgleich
\usepackage{setspace}					%Zeilenabstand: \onehalfspacing
\usepackage[utf8]{inputenc}				%Zeichenkodierung für Umlaute
\usepackage{scrlayer-scrpage}					%für Dokumentformatierung mit Pagestyle
\usepackage[version=4]{mhchem}						%Chemische Formeln (mit \ce )
\usepackage{chemfig}
\usepackage{siunitx}
\usepackage{caption}										%Überschriften formatieren
\captionsetup{labelfont={bf,it},tablename=Tab.,figurename=Abb.}				%Überschriften formatieren
\captionsetup[table]{singlelinecheck=false,justification=raggedright}
\captionsetup[figure]{justification=raggedright}				

%Abbildungs und Tabellennummerierung anpassen
\renewcommand*\figureformat {\figurename~\arabic{section}.\arabic{subsection}.\thefigure.}	%Punkt nach Abb. #
\renewcommand*\tableformat {\tablename~\arabic{section}.\arabic{subsection}.\thetable:}		%Doppelpunkt nach Tab. #
\renewcommand*\captionformat {~}								%Nur Abstand, kein Zeichen nach Tab. #: bzw Abb. #. 

%Literatur als Subsubsection definieren und von Literaturverzeichnis in Literatur umbenennen
\makeatletter
\renewcommand*\bib@heading{\subsubsection{\bibname}			
\@mkboth{\subsubsectionmarkformat \bibname}{\subsubsectionmarkformat \bibname}}
\makeatother
\addto\captionsngerman{\renewcommand*\bibname{Literatur}} 

\usepackage{prettyref}						%Nummerierung der Tabellen und Abbildungen im Text an Tab. x.x.# anpassen
%%% Für Abbildungen %%%
\newrefformat{fig}{Abb.~\arabic{section}.\arabic{subsection}.\ref{#1}}
%%% Für Tabellen %%%
\newrefformat{tab}{Tab.~\arabic{section}.\arabic{subsection}.\ref{#1}}

%Links im Dokument:
\usepackage[									%als letztes Paktet einbinden!
	linkcolor=black,							%Farbe der Links
	citecolor=black,							%Farbe der Link-Ziele
	urlcolor=black,								%Farbe von URLs
	colorlinks=true,							%Links nicht als Kasten, sondern direkt als Text anzeigen
]{hyperref}										%Links im PDF-File, unterstützung von URLs mit /url{}

%Dokumentformatierungen
\setkomafont{sectioning}{\bfseries} 	%Serifenschrift in Überschriften
\setkomafont{subsection}{\LARGE}
\setkomafont{subsubsection}{\large}
\pagenumbering{arabic}					%arabische Seitenzahlen
\geometry{left=3cm, right=2cm, top=2.5cm, bottom=2cm} 	%Seitenränder
\parindent0pt 									%Absätze einrücken um 0Pt
\parskip 10pt									%Abstand nach Absatz
\onehalfspacing								%Zeilenabstand
\pagestyle{scrheadings}					%Kopf- und Fussleistenformatierung
\headsep0.4cm								%Abstand Kopfleiste-Text
\footskip0.8cm									%Abstand Fussleiste-Text
\urlstyle{same}								%URLs normal formatieren, nicht als Typewriter

%Befehle definieren

%%%%% create chemequation environment

\newcounter{chemequation}
\renewcommand*\thechemequation{(\Roman{chemequation})}
\newenvironment{chemequation}{%
	\stepcounter{chemequation}%
	\begin{equation}}%
{\tag*{\thechemequation}%
	\end{equation}}

%%%%% create chemtag to iterate and ref chemequations in an align environment

\newcommand*\chemtag{%
	\stepcounter{chemequation}%
	\tag*{\thechemequation}}

%%%%% to create a chemicalequation float

\usepackage{float}
\newfloat{chemicalequation}{htbp}{cex}%[chapter]
\floatname{chemicalequation}{}

\newcommand{\zitat}[1]{\textsuperscript{\cite{#1}}}

\newcommand{\bild}[4]{					%Bilder mit Beschriftung einfügen
	\begin{figure}[#4]						%#4: Platzierung (htbp)
		\centering								%Bilder mittig einfügen
		\captionbox								%Beschriftungsumgebung
			{#2\label{fig: #1}}				%#2: Beschriftung
			{\includegraphics[#3]{#1}}	%#1: Dateiname, #3: Bildeigenschaften
	\end{figure}} 

\newcommand{\formel}[2]{				%Formel einfügen und Referenz erzeugen
	\begin{equation}
		\label{eq: #1}							%#1: Formelname
		#2											%Formelinhalt
	\end{equation}}

\newcommand{\cformel}[2][]{				%Chemische Formel einügen und Referenz erzeugen
	\begin{chemequation}
		\label{ceq: #1}					%#1: Formelname (optional)
		\ce{#2}							%#2: Formelinhalt
	\end{chemequation}\\}

	% 2. Ein neuer, eigener Befehl NUR für Formeln mit Zusatz auf der rechten Seite:
\newcommand{\cformelpot}[3]{%
    \begin{chemequation}
        \label{ceq: #1}
        \ce{#2} \hspace{2cm} #3
    \end{chemequation}%
}

\newcommand{\ox}[2]{
	\overset{\text{#2}}{\ce{#1}}}
	
\newcommand{\tabelle}[6]{				%Tabelle mit Beschriftung einfügen
\begin{table}[#3]							%#3: Platzierung (htbp)
\centering										%Tabelle mittig einfügen
\captionbox										%Beschriftungsumgebung
{#2												%#2: Tabellenbeschriftung
\label{tab: #1}}								%#1: Tabellenname
{\begin{tabular}{#4}						%#4: Anzahl und Ausrichtung Spalten (l c r)
\hline 											%horizontale Linie
#5													%#5: Obere Zeile
\hline
#6													%#6: Datenzeilen
\hline 
\end{tabular}}
\end{table}}

\usepackage{array,ragged2e}
\newcolumntype{P}[1]{>{\RaggedRight\arraybackslash}p{#1}} %Verhindert ungewollte Lücken in Tabellen

\newcommand{\langetabelle}[6]{  %Seitenübergreifende Tabelle, Kopfzeile wird 
	\begin{centering}           %auf jeder neuen Seite ebenfalls angezeigt
		{\begin{longtable}[#3]{#4}
				\caption{#2\label{tab: #1}}\\				
				\hline
				#5
				\hline
				\endfirsthead
                \caption{#2\label{tab: #1}}\\
				\hline
				#5
				\hline
				\endhead
				#6												
				\hline 
		\end{longtable}}
\end{centering}}

%Plots erstellen&&&&&&&&

%#1: Referenz
%#2: Name
%#3: x label
%#4: y label
%#5: x-werte,y-werte (Spaltenweise)
%#6: Quelle (.txt Datei)
%#7: "Nummerierung" (manuell) <.<

\newcommand{\plot}[6]{
	\begin{figure}[H]	
		\centering
		\begin{tikzpicture}
	
		\begin{axis}[
		title style={at={(0.5,0)},
		anchor=north,
		yshift=-40},
		xlabel={#3},
		ylabel={#4},
		legend pos=south east,
		legend entries={},
		]
		\addplot[color=black,mark=*,only marks] table [#5] {#6};
		\label{plot: plot1}
		\end{axis}
		\end{tikzpicture}
		\caption{#2}
		\label{fig: #1}
	\end{figure}
}

\newcommand{\molL}[1]{\ensuremath{\mathrm{\SI{#1}{\frac{mol}{L}}}}} % \molL{1} => \SI{1}{\frac{mol}{L}}
\newcommand{\mL}[1]{\ensuremath{\mathrm{\SI{#1}{mL}}}} % \mL{1} => \SI{1}{mL}

\newcommand{\Rn}[1]{\romannumeral#1\relax}                   %\Rn{1} = i
\newcommand{\RN}[1]{\uppercase\expandafter{\romannumeral#1}} %\RN{1} = I

%Dokument:%%%%%%%%%%%%%%%%%%%%%%%%%%%%%%%%%%%%%%%%%%%%%%%%%%%%%%%%%%%%%%%%%%%%%%%%%%%%
\begin{document}
	
	%Titelseite
	\begin{titlepage}
		\begin{center}
			\hrule													%waagerechte Linie
			\vspace*{5cm}
			\Huge {\textbf {Versuchsreihe 6}}\\						%Titel Versuchsreihe
			\vspace*{0.5cm}
			%		\large {\bf {Untertitel}}\\						%Untertitel
			\vspace*{1cm}
			\normalsize
			Protokoll	\\									
			Praktikum für Nanowissenschaftler:innen\\
			im Studiengang Nanowissenschaften\\
			\vspace*{0.5cm}
			\Large {\textbf{Michael Quint}}		\\						%Name
			\vspace*{0.5cm}
					\normalsize{Platz 9a}								%Platznummer
			\vfill
		\end{center}
		\begin{tabbing}
			\hspace{3cm}\=\kill
			Assistenten: \> Lisa Huber, Maximilian Ruffer \\						% Name Assistent
			Version: \> 1 \\									%Version
			Datum: \> 08.11.2026 \\								%Datum
		\end{tabbing}
		\hrule														%Waagerechte Linie
\end{titlepage}

\setcounter{page}{1}
\setcounter{section}{6}								%Versuchstag #.x.x bzw Versuchsreihe
\setcounter{subsection}{2}							%Versuchsnummer x.#.x startet bei 0 für 1 
\setcounter{secnumdepth}{4}
	
\subsection{Titration einer starken Säure}

\subsubsection{Einleitung und Aufgabenstellung}

Es war die Konzentration einer Salzsäurelösung mittels Titration in zwei Druchgängen zu bestimmen.
Dabei waren die gemessenen pH-Werte aus dem ersten Durchgang graphisch aufzutragen.

\subsubsection{Theoretischer Hintergrund}

Die erste allgemeingültige Säure-Base-Definition wurde von Arrhenius 1883 entwickelt.
Nach dieser Definition sind Säuren (\ce{HA}) Stoffe, die unter Bildung von Hydronium-Ionen (\ce{H3O+}) dissoziieren.
\cformel[s1]{HA + H2O ->[][Dissoziation] H3O+ + A-}
Basen (\ce{XOH}) sind darauffolgend Stoffe, die unter Bildung von Hydroxid-Ionen (\ce{OH-}) dissoziieren.\zitat{Arr}
\cformel[c1]{XOH ->[][Dissoziation] X+ + OH-}

Da die Definition von Arrhenius streng genommen nur für wässrige Lösungen gilt, wurde eine allgemeine Erweitung der Säure-Base-Definition 1923 von Br\o nsted entwickelt.
Nach seiner Theorie sind Säuren Protonendonatoren, also jene Moleküle bzw. Ionen, die in der Lage sind, Protonen (\ce{H+}-Ionen) abzuspalten.
Eine Protonenübertragungsreaktion nennt man auch Protolyse.
Basen sind korrespondierend dazu Protonenakzeptoren, also Moleküle bzw. Ionen, welche in der Lage sind, Protonen aufzunehmen.\\
Zusammengefasst sind also bei Protolysereaktionen zwangsweise immer zwei Säure-Base-Paare beteiligt,
in welchen jeder Stoff nach einer Protonendonation bzw. -akzeptanz zu einer konjugierten Base bzw. Säure reagiert.\zitat{konj}
\cformel[sb]{\underset{\text{Säure}}{\ce{HA}} + \underset{\text{Base}}{\ce{B}} <=> \underset{\text{konj. Base}}{\ce{A-}} + \underset{\text{konj. Säure}}{\ce{HB+}}}
Wasser (\ce{H2O}) hat als sog. Ampholyt außerdem die Eigenschaft, sowohl als Säure als auch als Base zu reagieren.\zitat{konj}
\cformel[sb2]{\underset{\text{Säure}}{\ce{H2O}} + \underset{\text{Base}}{\ce{H2O}} <=> \underset{\text{konj. Base}}{\ce{OH-}} + \underset{\text{konj. Säure}}{\ce{H3O+}}}

Protolysereaktionen sind wie alle anderen chemischen Reaktionen ebenfalls Gleichgewichtsreaktion.
Im Falle von Wasser spricht man auch von einer Autoprotolyse.
\cformel[mwg1]{2H2O_{(l)} <=> H3O^{+}_{(aq)} + OH^{-}_{(aq)}}
Das Massenwirkungsgesetz (MWG) lässt sich demnach wie folgt schreiben, wobei $a$ die Aktivität eines Stoffes beschreibt:
\formel{mwg2}{K = \frac{a(\ce{H3O+})\cdot a(\ce{OH-})}{a(\ce{H2O})^2}}
Da die Aktivität von Wasser nahezu konstant ist, wird das Ionenprodukt des Wassers $K_{\mathrm{W}}$ eingeführt:
\formel{mwg3}{K \cdot a(\ce{H2O})^2 = K_{\mathrm{W}} = a(\ce{H3O+})\cdot a(\ce{OH-})}
Bei $T$ = 298 K beträgt $K_{\mathrm{W}}$ = $1,0 \cdot 10^{-14}$. Bei einer negativen dekadischen Logarithmierung ($\log_{10}x \equiv \lg x$) ergibt sich folgende Beziehung:
\formel{mwg4}{-\lg(K_{\mathrm{W}}) = -\lg(a(\ce{H3O+})\cdot a(\ce{OH-})) = -\lg(a(\ce{H3O+})) -\lg(a(\ce{OH-}))}
Mit der Einführung der pH- und pOH-Werte gelten demnach folgende Beziehungen\zitat{MWG}:
\formel{pH}{\mathrm{pH} \equiv -\lg(a(\ce{H3O+}))}
\formel{pOH}{\mathrm{pOH} \equiv -\lg(a(\ce{OH-}))}
\formel{fin}{\mathrm{pH} + \mathrm{pOH} = -\lg(1 \cdot 10^{-14}) = 14}




\subsubsection{Experimenteller Teil}

\bild
		{pH_stark.png}					%Dateiname des Bildes, eindeutiger Name des Bildes, Verweis: \prettyref{fig: Bildname}
		{Gemessene pH-Werte aufgetragen gegen das Volumen der zugegebenen NaOH-Maßlösung.%Text, der unter dem Bild erscheint
		}			%Name der Quelle im Literaturverzeichnis
		{scale=1}		%angle, draft, scale, height, width
		{h}				%h: here, t: top, b: bottom, p: page; bestimmt, wo das Bild eingefügt wird. p: spezielle Seite für Abbildungen und Tabellen, Kombinationen möglich, Versuch der Platzierung gemäß Reihenfolge.

\subsubsection{Ergebnisse und Diskussion}

\subsubsection{Zusammenfassung}

\subsubsection{Sicherheit und Entsorgung}
		
Die für den Versuch verwendeten Chemikalien sind in \prettyref{tab: Chemikalien} mit Gefahrensymbolen, H- und P-Sätzen sowie Hinweisen zur Entsorgung aufgeführt.\zitat{test}\zitat{indikator1} Verwendete KMR-Stoffe sind in \prettyref{tab: KMR_leer} aufgeführt.\zitat{test}
	
\langetabelle
	{Chemikalien}
	{Verwendete Chemikalien.}
	{H}  % hier ggf. ein großes H eintragen, wenn die Tabelle sonst im Anhang verschwindet
	{l p{1.9cm} p{2.2cm} p{4.5cm} c}
	{Substanz & Symbol & H-Sätze & P-Sätze & Entsorgung \\}
	{ % Erster Stoff

	\ce{NaOH}
		& \begin{tabular}[t]{l} GHS05 \end{tabular}
		& 290, 314
		& 234, 260, 280, \newline 303+361+353, \newline 304+340+310, \newline 305+351+338,
		& (1) \\

	\ce{HCl} (konz.)
		& \begin{tabular}[t]{l} GHS05 \\ GHS07 \end{tabular}
		& 290, 314 \newline 335 
		& 234, 261, 271, 280, \newline 303+361+353, \newline 305+351+338,
		& (2) \\

	\ce{Tymolphthalein 0,1 \% in Ethanol}
		& \begin{tabular}[t]{l} GHS02 \\ GHS07 \end{tabular}
		& 225, 319 
		&  210, 233, 240, 241, \newline 242, 305+351+338
		& (3) \\
	}

\langetabelle
	{KMR_leer}
	{Verwendete kanzerogene, mutagene oder reproduktionstoxische Gefahrstoffe (KMR).}
	{H}
	{l P{3.8cm} P{2.4cm} P{3.5cm} c}
	{Substanz & H-Sätze & KMR-Einstufung & Verwendete Menge & Entsorgung \\}
	{ - & - & - & - & - \\}
	
Entsorgung: \\
\begin{itemize}
	\item[(1)] Die Lösung wurde im Behälter für andere Basen entsorgt.
	\item[(2)] Die Lösung wurde im Behälter für nicht-nitrathaltige Säuren entsorgt.
    \item[(3)] Die titrierte Lösung wurde im Behälter für andere Basen entsorgt.
\end{itemize}	


\begin{thebibliography}{4} %manuelles Quellenverzeichnis

	\bibitem{Arr}
						{E. Riedel, C. Janiak}, 
			\textit	{Anorganische Chemie}, 
						{10. Auflage}, 
						{de Gruyter}, 
						{Berlin}, 
			\textbf	{2022}, 
			S.			{345-346}.

	\bibitem{konj}
						{E. Riedel, C. Janiak}, 
			\textit	{Anorganische Chemie}, 
						{10. Auflage}, 
						{de Gruyter}, 
						{Berlin}, 
			\textbf	{2022}, 
			S.			{346-347}.

	\bibitem{MWG}
						{E. Riedel, C. Janiak}, 
			\textit	{Anorganische Chemie}, 
						{10. Auflage}, 
						{de Gruyter}, 
						{Berlin}, 
			\textbf	{2022}, 
			S.			{348-352}.

	\bibitem{test}
						{DGUV}, 
						„{GESTIS-Stoffdatenbank}“, zu finden unter 
			\url		{https://gestis.dguv.de/}, 
			\textbf	{01.11.2026}.
	
	\bibitem{indikator1}
						{Carl Roth}, 
						„{Sicherheitsdatenblatt: Thymolphthalein - Indikatorlösung}“, zu finden unter 
			\url		{https://www.carlroth.com/}, 
			\textbf	{01.11.2026}.
			
\end{thebibliography}

%---------------------------------------------------------------------------------------------------------------------------------
\newpage

\setcounter{subsection}{3}							%Versuchsnummer x.#.x startet bei 0 für 1 
\setcounter{secnumdepth}{4}
	
\subsection{Titration einer schwachen Säure}

\subsubsection{Einleitung und Aufgabenstellung}


\subsubsection{Theoretischer Hintergrund}


\subsubsection{Experimenteller Teil}


\subsubsection{Ergebnisse und Diskussion}

\subsubsection{Zusammenfassung}

\subsubsection{Sicherheit und Entsorgung}

Die für den Versuch verwendeten Chemikalien sind in \prettyref{tab: Chemikalien3} mit Gefahrensymbolen, H- und P-Sätzen sowie Hinweisen zur Entsorgung aufgeführt.\zitat{hauke3}\zitat{indikator2} Verwendete KMR-Stoffe sind in \prettyref{tab: KMR_leer3} aufgeführt.\zitat{hauke3}		

\langetabelle
	{Chemikalien2}
	{Verwendete Chemikalien.}
	{H}  % hier ggf. ein großes H eintragen, wenn die Tabelle sonst im Anhang verschwindet
	{l p{1.9cm} p{2.2cm} p{4.5cm} c}
	{Substanz & Symbol & H-Sätze & P-Sätze & Entsorgung \\}
	{ % Erster Stoff

	\ce{NaOH}
		& \begin{tabular}[t]{l} GHS05 \end{tabular}
		& 290, 314
		& 234, 260, 280, \newline 303+361+353, \newline 304+340+310, \newline 305+351+338,
		& (1) \\

	\ce{CH3COOH}
		& \begin{tabular}[t]{l} GHS02 \\ GHS05 \end{tabular}
		& 226, 314 
		& 210, 280 \newline 301+330+331, \newline 303+361+353, \newline 305+351+338, 310
		& (2) \\


	\ce{Tymolphthalein 0,1 \% in Ethanol}
		& \begin{tabular}[t]{l} GHS02 \\ GHS07 \end{tabular}
		& 225, 319 
		&  210, 233, 240, 241, \newline 242, 305+351+338
		& (3) \\

    }

\langetabelle
	{KMR_leer2}
	{Verwendete kanzerogene, mutagene oder reproduktionstoxische Gefahrstoffe (KMR).}
	{H}
	{l P{3.8cm} P{2.4cm} P{3.5cm} c}
	{Substanz & H-Sätze & KMR-Einstufung & Verwendete Menge & Entsorgung \\}
	{ - & - & - & - & - \\}
	
Entsorgung: \\
\begin{itemize}
    \item[(1)] Die Lösung wurde im Behälter für andere Basen entsorgt.
	\item[(2)] Die Lösung wurde im Behälter für nicht-nitrathaltige Säuren entsorgt.
    \item[(3)] Die titrierte Lösung wurde im Behälter für andere Basen entsorgt.
\end{itemize}	

\begin{thebibliography}{4}		%manuelles Quellenverzeichnis

	
	\bibitem{hauke2}
						{DGUV}, 
						„{GESTIS-Stoffdatenbank}“, zu finden unter 
			\url		{https://gestis.dguv.de/}, 
			\textbf	{01.11.2026}.

	\bibitem{indikator2 }
						{Carl Roth}, 
						„{Sicherheitsdatenblatt: Thymolphthalein - Indikatorlösung}“, zu finden unter 
			\url		{https://www.carlroth.com/}, 
			\textbf	{01.11.2026}.


\end{thebibliography}

%--------------------------------------------------------------------------------------------------------------------------------------
\newpage

\setcounter{subsection}{4}							%Versuchsnummer x.#.x startet bei 0 für 1 
\setcounter{secnumdepth}{4}
	
\subsection{Pufferwirkung}

\subsubsection{Einleitung und Aufgabenstellung}


\subsubsection{Theoretischer Hintergrund}


\subsubsection{Experimenteller Teil}


\subsubsection{Ergebnisse und Diskussion}

\subsubsection{Zusammenfassung}

\subsubsection{Sicherheit und Entsorgung}
		
Die für den Versuch verwendeten Chemikalien sind in \prettyref{tab: Chemikalien3} mit Gefahrensymbolen, H- und P-Sätzen sowie Hinweisen zur Entsorgung aufgeführt.\zitat{hauke3}\zitat{indikator3} Verwendete KMR-Stoffe sind in \prettyref{tab: KMR_leer3} aufgeführt.\zitat{hauke3}
	
\langetabelle
	{Chemikalien3}
	{Verwendete Chemikalien.}
	{H}  % hier ggf. ein großes H eintragen, wenn die Tabelle sonst im Anhang verschwindet
	{l p{1.9cm} p{2.2cm} p{4.5cm} c}
	{Substanz & Symbol & H-Sätze & P-Sätze & Entsorgung \\}
	{ % Erster Stoff

	\ce{NaOH}
		& \begin{tabular}[t]{l} GHS05 \end{tabular}
		& 290, 314
		& 234, 260, 280, \newline 303+361+353, \newline 304+340+310, \newline 305+351+338,
		& (1) \\

	\ce{CH3COOH}
		& \begin{tabular}[t]{l} GHS02 \\ GHS05 \end{tabular}
		& 226, 314 
		& 210, 280 \newline 301+330+331, \newline 303+361+353, \newline 305+351+338, 310
		& (2) \\

	\ce{HCl} (1 mol/L)
		& \begin{tabular}[t]{l} GHS05 \end{tabular}
		& 290
		& 234, 390
		& (2) \\

	\ce{Tymolphthalein 0,1 \% in Ethanol}
		& \begin{tabular}[t]{l} GHS02 \\ GHS07 \end{tabular}
		& 225, 319 
		&  210, 233, 240, 241, \newline 242, 305+351+338
		& (3) \\

	\ce{Tymolblau in Ethanol 0,05 \%}
		& \begin{tabular}[t]{l} GHS02 \\ GHS07 \end{tabular}
		& 225, 319
		& 210, 280, \\ 303+361+353, \\ 305+351+338
		& (3) \\

	\ce{NH3} (25 \%)
		& \begin{tabular}[t]{l} GHS05 \\ GHS07 \\ GHS09 \end{tabular}
		& 314, 335, 410
		& 261, 271, 273, 280, \newline 303+361+353, \newline 305+351+338
		& (4) \\
    }

\langetabelle
	{KMR_leer3}
	{Verwendete kanzerogene, mutagene oder reproduktionstoxische Gefahrstoffe (KMR).}
	{H}
	{l P{3.8cm} P{2.4cm} P{3.5cm} c}
	{Substanz & H-Sätze & KMR-Einstufung & Verwendete Menge & Entsorgung \\}
	{ - & - & - & - & - \\}
	
Entsorgung: \\
\begin{itemize}
    \item[(1)] Die Lösung wurde im Behälter für andere Basen entsorgt.
	\item[(2)] Die Lösung wurde im Behälter für nicht-nitrathaltige Säuren entsorgt.
    \item[(3)] Die titrierte Lösung wurde im Behälter für andere Basen entsorgt.
    \item[(4)] Die Lösung wurde mit Salzsäure angesäuert und im Behälter für nicht-nitrathaltige Säuren entsorgt.
\end{itemize}	


\begin{thebibliography}{4}		%manuelles Quellenverzeichnis

	\bibitem{hauke3}
						{DGUV}, 
						„{GESTIS-Stoffdatenbank}“, zu finden unter 
			\url		{https://gestis.dguv.de/}, 
			\textbf	{01.11.2026}.
	
	\bibitem{indikator3}
						{Carl Roth}, 
						„{Sicherheitsdatenblatt: Thymolblau Indikatorlösung}“, zu finden unter 
			\url		{https://www.carlroth.com/}, 
			\textbf	{01.11.2026}.


\end{thebibliography}

\end{document}

%\newpage

%%%%%%%%%%%%%%%%%%%%%%%%%%%%%%%%%%%%%%%%%%%%%%%%%%%%%%%%%%%%%%%
Neuer Versuch:
\subsection{Versuchsname}

\subsubsection{Einleitung und Aufgabenstellung}
\subsubsection{Theoretischer Hintergrund}
\subsubsection{Experimenteller Teil}
\subsubsection{Ergebnisse und Diskussion}
\subsubsection{Zusammenfassung}
\subsubsection{Sicherheit und Entsorgung}	
Die für den Versuch verwendeten Chemikalien sind in \prettyref{tab: entsorgung1} mit Gefahrensymbolen, H- und P-Sätzen sowie Hinweisen zur Entsorgung aufgeführt.\zitat{Gefahren1} Verwendete KMR-Stoffe sind in \prettyref{tab: KMR1} aufgeführt.
	
\langetabelle
	{entsorgung1}
	{Verwendete Chemikalien.}
	{H}  % hier ggf. ein großes H eintragen, wenn die Tabelle sonst im Anhang verschwindet
	{l p{1.4cm} p{2.2cm} p{4.5cm} c}
	{Substanz & Symbol & H-Sätze & P-Sätze & Entsorgung \\}
	{\\}
    

\langetabelle
	{KMR1}
	{Verwendete kanzerogene, mutagene oder reproduktionstoxische Gefahrstoffe (KMR).}
	{H}
	{l P{3.8cm} P{2.4cm} P{3.5cm} c}
	{Substanz & H-Sätze & KMR-Einstufung & Verwendete Menge & Entsorgung \\}
	{ - & - & - & - & - \\}
Entsorgung: \\
\begin{itemize}
    \item[(1) Die verwendete Lösung wurde mit Salzsäure angesäuert und im Behälter für nicht-nitrathaltige Säuren entsorgt.] 
\end{itemize}	

\newpage

\begin{thebibliography}{4}		%manuelles Quellenverzeichnis
	\bibitem{Gefahren1}
						{H. Heller}, 
			\textit	{Grundlagenpraktikum der Chemie für Studierende der Nanowissenschaften}, 
						{Version 1.18},
                        {Universität Hamburg},
						{Hamburg}, 
			\textbf	{2023}, 
			S.{34}.
			
			
\end{thebibliography} 
\newpage

Nützliche Beispiele:

%Tabelle für Reagenzgläserinhalte
\tabelle
	{Tabellenname}												%Eindeutige Bezeichnung der Tabelle. \prettyref{tab: Tabellenname}
		{Tabellenüberschrift.}	%Inhalt der Reagenzgläser ... bis...
		{htb}	
		{l l l}					
		{Reagenzglas & Inhalt & Beschreibung}				
		{1 & 2~mL NaOH & klare, farblose Lsg.\\
		2 & & 
		}										


%Formel Löslichkeit KClO4 -> chemische Formel
\begin{align}
    \chemtag
    \ce{KClO4(aq) <=> K+(aq) + ClO4-(aq)}
\end{align}

%hier allgemeine Formel 
\formel
		{formel2}				%eindeutige Bezeichnung der Formel, Verweis: \ref{eq: formelname}
		{L_{\mathrm{A}_{a} B_{b}} = \sqrt[m+n]{\frac{K_{\mathrm{L}}}{a^{a} \cdot b^{b}}}  }


Bilder:
	\bild
		{bildname}					%Dateiname des Bildes, eindeutiger Name des Bildes, Verweis: \prettyref{fig: Bildname}
		{Bildbeschreibung.		%Text, der unter dem Bild erscheint
		\zitat{Quelle}}			%Name der Quelle im Literaturverzeichnis
		{Bildeigenschaften}		%angle, draft, scale, height, width
		{Platzierung}				%h: here, t: top, b: bottom, p: page; bestimmt, wo das Bild eingefügt wird. p: spezielle Seite für Abbildungen und Tabellen, Kombinationen möglich, Versuch der Platzierung gemäß Reihenfolge.
	
		
	Bildeigenschaften:
		angle - Bestimmt den Drehwinkel um einen Referenzpunkt für die Grafik. Positive Werte drehen die Grafik gegen den Uhrzeigersinn, negative mit ihm. 
		draft - Ist draft gesetzt, wird nur ein Platzhalter der Grafik angezeigt. Beschleunigt den Übersetzungsvorgang des Dokuments erheblich. 
		scale - Skaliert die Grafik anhand eines Skalierungsfaktors. 
		height - Skaliert die Grafik auf die angegebene Höhe. Die Angabe der Größe muss in einer LaTeX-spezifischen Einheit erfolgen. 
		width - Skaliert die Grafik auf die angegebene Breite. Die Angabe der Größe muss in einer LaTeX-spezifischen Einheit erfolgen. 	
	

Formeln:
	\formel
		{formelname}				%eindeutige Bezeichnung der Formel, Verweis: \ref{eq: formelname}
		{Formel}						%


Tabellen:
	\tabelle
		{Tabellenname}												%Eindeutige Bezeichnung der Tabelle. \prettyref{tab: Tabellenname}
		{Tabellenüberschrift.\zitat{Quelle}}						%Text, der über der Tabelle erscheint, Quelle der Daten
		{Platzierung}													%h: here, t: top, b: bottom, p: page; bestimmt, wo das Bild eingefügt wird. p: spezielle Seite für Abbildungen und Tabellen, Kombinationen möglich, Versuch der Platzierung gemäß Reihenfolge.		
		{Spaltenanzahl_und_Ausrichtung}						%{ l c r }: 3 Spalten, links mittig und rechts; { c c c c }: 4 Spalten, alle mittig
		{Titelzeile}														%Obere Zeile(n), danach folgt ein waagerechter Strich, Spalten durch & trennen, Neue Zeile durch \\
		{Datenzeilen}													%Zeilen unter dem waagerechten Strich, Spalten durch & trennen, Neue Zeile durch \\
		
	Zellen verbinden:	
		\multicolumn{Anzahl_der_verbunden_Spalten}{Ausrichtung (l/c/r}{Inhalt}
	
Reaktionen:
Formeln: 
\begin{center}
\ce{NH+4}, \ce{Fe3(CO)12}, \ce{Al^{3+}},\ce{Cr2O^{2-}7}
\end{center}

Isotope:
\begin{center}
\ce{^{37}_{17}Cl}, \ce{^{32}_{16}S^{2-}}
\end{center}

Bindungen: 
\begin{center}
\ce{H-O-N=S}, \ce{Cl-C#N}
\end{center}

Reaktionsgleichungen: 
\begin{center}
\ce{2H2 + O2 -> 2H2O}
\ce{MgCl2 ->[\text{Elektrolyse}] Mg\,(l) + Cl_{2}\,(g)}
\end{center}

Redoxreaktion:
\begin{center}
$\overset{0}{\ce{S}} + \overset{0}{\ce{O2}} \ce{->} \overset{+\text{IV}}{\ce{S}}+\overset{-\text{II}}{\ce{O2}}$
\end{center}


Literaturverweise:
		\bibitem{zeitschrift}
						{A. Autor}, 
			\textit	{Zeitschrift} 
			\textbf	{Jahr}, 
			\textit	{Band} 
						{Seite-Seite}.
		
		\bibitem{buch_ohne_herausgeber}
						{A. Autor}, 
			\textit	{Titel}, 
						{Verlag}, 
						{Ort}, 
			\textbf	{Jahr}, 
			S.			{Seite-Seite}.
			
		\bibitem{buch_mit_herausgeber}
						{A. Autor}, 
			\textit	{Titel}, 
						{Auflage}, 
						{Verlag}, 
						{Ort}, 
			\textbf	{Jahr}, 
			S.			{Seite-Seite}.
			
		\bibitem{dissertation}
						{A. Autor}, 
						{Dissertation}, 
						{Universität}, 
			\textbf	{Jahr}.
			
		\bibitem{internetseite}
						{A. Autor}, 
						„{Titel}“, zu finden unter 
			\url		{www.xyz.de}, 
			\textbf	{Jahr}.

% Kommentare Sophie:

% <- Für Kommentare oder Code der gerade nicht gebraucht wird, wird beim kompilieren nicht beachtet
%Um Bilder einzufügen: \bild{'Dateiname' (Datei muss im gleichem Ordner mit der TeX Datei liegen)}{'Bildunterschrift'}{width= ...cm (Größe des Bildes)}{H (Ort des Bildes, H=Hier)}

%Für Summenformeln im Text: \ce{z.B.: H2SO4 (Zahlen werden automtisch tief gestellt)}

% Um Formeln im Text zu referieren: \prettyref{fig (für Bilder und Graphen)/ eq (für mathematische Formeln)/ ceq (für chemische Formeln): 'Referenz'}

%Für chemische Formeln: \cformel['Referenz']{2\ox{Fe}{+\RN{3}}Cl3 + \ox{Fe}{+\RN{2}}Cl2 + 8NH4+ + 8 OH- \ce{<=>[{N2H4}]} \ox{Fe}{+\RN{2}}O\cdot \ox{Fe}{+\RN{3}}2O3 \ce{v} + 4 H2O + 8NH4+ + 4 Cl-}

%Für mathematische Formeln: \formel{'Referenz'}{2\cdot d\cdot \ce{sin}(\theta)= n \lambda}

%Tabellen: Online Converter um Exel Tabellen in LaTeX Format zu bringen: https://www.tablesgenerator.com
%\tabelle{'Referenz'}{'Tabellenüberschrift'}{H}{'Breite der Spalte bzw. Textausrichtung' (Spaltenbreite passt sich Textlänge an: l=links, r=rechts, c=zentriert. Spaltenbreite fest und Blocksatz: p{...cm}), Jede Spalte muss einzeln definiert werden, z.B l l l = 3 Spalten}{'Spaltenüberschrift' muss der Anzahl an Spalten entsprechen. Spaltenüberschriften müssen durch ein & getrennt werden. Außerdem wichtig am Ende: \\}{'Tabelleneinträge', müssen ebenfalls durch ein & getrennt werden und am Ende der Zeile ein \\ haben
%	- & - & - & \\
%	- & - & - & \\}
%Beispiel für eine Tabelle (zum Ausführen '%' entfernen):
%\tabelle{Auflösungsvermögen}{Experimentell ermitteltes und theoretisches Auflösungsvermögen für Objektiv 20X und 63X bei schräger und gerader Beleuchtung.}{H}{l l p{2cm} p{3cm}  p{3cm}}{Objektiv&Beleuchtung& Spalt mit einem Peak&experimentelles Auflösungsvermögen~/~nm&theoretisches Auflösungsvermögen~/~nm\\}
%	{20X&gerade&16&>750&1125\\
%	20X&schräg&9&690&563\\
%	63X&gerade&11&>750&600\\
%	63X&schräg&7&584&300\\}